%%%%%%%% Paper: simple cover-time prediction on small graphs %%%%%%%%%%%%%%%%%
\documentclass{article}
\usepackage{microtype}
\usepackage{graphicx}
\usepackage{booktabs}
\usepackage{hyperref}
\newcommand{\theHalgorithm}{\arabic{algorithm}}
\usepackage{icml2025}
\usepackage{amsmath,amssymb}
\usepackage[capitalize,noabbrev]{cleveref}

\icmltitlerunning{Simple statistics predict cover times on small graphs}

\begin{document}

\twocolumn[
\icmltitle{Simple Graph Statistics Predict Random-Walk Cover Times\\on Small Graphs (n$\le$30)}

\icmlsetsymbol{equal}{*}
\begin{icmlauthorlist}
\icmlauthor{Anonymous Author}{anon}
\end{icmlauthorlist}
\icmlaffiliation{anon}{Anonymous Institution}
\icmlcorrespondingauthor{Anonymous Author}{anon@example.org}

\icmlkeywords{random walks, cover time, Markov chains, empirical prediction}

\vskip 0.3in
]

\printAffiliationsAndNotice{}

\begin{abstract}
The cover time of a simple random walk is the expected number of steps needed to visit every vertex. Classical bounds such as $2m(n-1)$ are tight only in the worst case and can overestimate the cover time of a small graph by an order of magnitude. We ask whether four easily computed statistics (vertex count $n$, edge count $m$, maximum degree $\Delta$, diameter $D$) give a more accurate prediction. Using 834 connected graphs with $3\le n\le 30$ from 15 synthetic families, and Monte Carlo estimates with 300 walks per start vertex, we fit a log-linear model $\log C \approx \beta_0 + \beta_1\log n + \beta_2\log m + \beta_3\log\Delta + \beta_4\log D$. Its median relative error under 5-fold cross-validation is $0.105$ (95\% bootstrap CI $[0.096, 0.114]$), against $8.5$ for $2m(n-1)$ and $0.50$ for the Matthews bound. Under leave-one-family-out evaluation the median error rises to $0.138$, and the model fails most on lollipop, grid, and complete graphs, where cover time depends on structure that these four statistics do not capture.
\end{abstract}

\section{Introduction}
\label{sec:intro}
For a connected graph $G$ with $n$ vertices and $m$ edges, the simple random walk started at $v$ has cover time $C_v$, and $C(G)=\max_v C_v$. \citet{aleliunas1979random} showed $C(G)\le 2m(n-1)$, and \citet{feige1995tight} showed the worst case over graphs is $\Theta(n^3)$, with the lollipop family attaining the leading constant. Survey treatments appear in \citet{lovasz1993random} and \citet{aldous2014reversible}. These results describe the extremal graph, not a typical small one. A user who needs the expected cover time of a specific 20-vertex network, for example to size a sampling budget or to compare traversal strategies, gets an answer that can be off by a factor of ten.

We ask a narrow empirical question: for connected graphs on at most 30 vertices, can a handful of cheap statistics predict $C(G)$ more accurately than the generic bounds, and where does such a predictor break down? Our contribution is (i) a benchmark of 834 small graphs with Monte Carlo cover times and their uncertainty, (ii) a four-feature log-linear predictor with bootstrap-quantified accuracy and an explicit leave-one-family-out evaluation, and (iii) a per-family error analysis identifying where the predictor fails. We do not claim a new theorem.

\section{Method}
\label{sec:method}
\paragraph{Graphs.} We generate 834 connected graphs with $3\le n\le 30$ from 15 families: paths, cycles, stars, wheels, complete graphs ($n\le 20$), complete bipartite graphs, grids, binary trees, lollipops, barbells, random labeled trees, Erd\H{o}s--R\'enyi $G(n,p)$ conditioned on connectivity, random 3-regular graphs, Watts--Strogatz small-world graphs, and Barab\'asi--Albert graphs. Random families use six seeds per size. All generators are in \texttt{src/graphs.py}; graphs are built with NetworkX \citep{hagberg2008exploring}.

\paragraph{Cover-time ground truth.} For each graph we simulate $300$ walks from each start vertex (all vertices if $n\le 12$, otherwise 10 uniformly chosen starts) and record the first time all vertices are visited. Our target is $\hat C(G)=\max_v \bar C_v$, the largest sample mean over starts. The max over noisy means is slightly biased upward; with 300 walks the per-start standard error is 2--5\% (median relative standard error per graph is reported in \cref{tab:families}). As a check, for the 140 graphs with $n\le 8$ we compute the exact cover time by solving the linear system on (current vertex, visited set) states. The Monte Carlo value agrees with the exact value to a median of $2.4\%$ and at most $12\%$, which is in line with the stated sampling noise.

\paragraph{Predictors.} We compare four predictors. (a) $2m(n-1)$, the bound of \citet{aleliunas1979random}. (b) $\kappa\cdot 2m(n-1)$ with one fitted constant. (c) A log-linear model on $(\log n, \log m)$. (d) A log-linear model on $(\log n,\log m,\log\Delta,\log D)$, fitted by ordinary least squares on $\log \hat C$. Predictor (d) is the main proposal. For reference we also compute the Matthews bound \citep{matthews1988covering} $H_{n-1}\max_{u,v}\mathbb E_u[T_v]$, which requires all hitting times and is therefore not a simple statistic.

\paragraph{Evaluation.} We report the median relative error $|\hat C_{\rm pred}-\hat C|/\hat C$ and the fraction of graphs predicted within 20\%. We use two protocols: repeated 5-fold cross-validation on graphs (5 repetitions), and leave-one-family-out (LOFO). Confidence intervals are 2000-sample bootstraps over graphs of the median error; they resample the per-graph errors and do not refit the model, so they capture sampling variability over the graph set but not the Monte Carlo noise in $\hat C$.

\section{Experiments}
\label{sec:exp}
Everything runs on a single CPU in about two minutes (\texttt{run.py}, \texttt{analyze.py}; seed 0). \cref{tab:main} summarises the four predictors, and \cref{fig:scatter} plots predictions against measurements.

\begin{table}[t]
\caption{Prediction accuracy over 834 graphs. Relative error is $|{\rm pred}-\hat C|/\hat C$. CI is a 95\% bootstrap over graphs for the CV median.}
\label{tab:main}
\begin{center}
\begin{small}
\begin{tabular}{lcccc}
\toprule
Predictor & CV median & 95\% CI & LOFO median & within 20\% (LOFO)\\
\midrule
$2m(n-1)$ (bound) & 8.49 & [7.89, 9.22] & 8.49 & 0\% \\
$\kappa\cdot 2m(n-1)$ & 0.665 & [0.59, 0.72] & 0.748 & 16\% \\
log-linear $(n,m)$ & 0.200 & [0.19, 0.22] & 0.276 & 37\% \\
\textbf{log-linear $(n,m,\Delta,D)$} & \textbf{0.105} & [0.096, 0.114] & \textbf{0.138} & \textbf{67\%} \\
Matthews (oracle) & 0.504 & -- & -- & 35\% (CV) \\
\bottomrule
\end{tabular}
\end{small}
\end{center}
\vskip -0.1in
\end{table}

\paragraph{Main result.} The four-feature log-linear model fits $\log\hat C$ with $R^2=0.965$ on the log scale. Its coefficients (intercept $-0.38$; $\log n$: $0.85$ [0.75, 0.98]; $\log m$: $-0.075$ [$-0.12$, $-0.04$]; $\log\Delta$: $0.68$ [0.63, 0.73]; $\log D$: $1.24$ [1.16, 1.30]) are bootstrap intervals over graphs. The sign of $\log m$ is negative once $n$, $\Delta$ and $D$ are controlled: at fixed diameter and maximum degree, denser graphs are covered faster. The diameter is the largest single contributor, consistent with a cover time that scales with how far the walk must travel. The predictor beats the generic bound by roughly a factor of $10$ in median error and beats the oracle Matthews bound on median error, although the Matthews bound is the tighter of the two in the families where hitting times are nearly uniform (star, complete).

\paragraph{Where it fails.} \cref{tab:families} gives LOFO median errors per family. The model is most accurate on wheels, bipartite graphs, and paths and cycles. It is least accurate on lollipops (0.52), grids (0.43), complete graphs (0.42), and binary trees (0.38). The lollipop result is the most informative. Its cover time is dominated by the time to escape a dense clique and then traverse a path, which is a bottleneck quantity that a single diameter does not describe. The complete-graph result is partly an extrapolation artefact: complete graphs are the only family with $D=1$, so leaving them out removes all training data at that diameter. We therefore read it as a warning about out-of-range extrapolation rather than as evidence about structure.

\begin{table}[t]
\caption{Leave-one-family-out median relative error of the main predictor, with the $2m(n-1)$ bound and the Matthews bound for comparison. Full table in \texttt{data/per\_family.csv}.}
\label{tab:families}
\begin{center}
\begin{small}
\begin{tabular}{lrrrr}
\toprule
Family & \#graphs & log-lin & bound & Matthews\\
\midrule
Lollipop & 27 & 0.52 & 0.93 & 2.10\\
Grid & 28 & 0.43 & 4.67 & 0.96\\
Complete & 18 & 0.42 & 37.99 & 0.03\\
Binary tree & 12 & 0.38 & 1.21 & 0.55\\
Random $G(n,p)$ & 84 & 0.17 & 16.19 & 0.65\\
Star & 28 & 0.16 & 3.61 & 0.01\\
Random tree & 84 & 0.13 & 1.08 & 1.09\\
Path & 28 & 0.11 & 0.54 & 1.61\\
Bipartite & 196 & 0.07 & 30.98 & 0.05\\
Wheel & 27 & 0.07 & 10.79 & 0.10\\
\bottomrule
\end{tabular}
\end{small}
\end{center}
\vskip -0.1in
\end{table}

\begin{figure}[t]
\begin{center}
\centerline{\includegraphics[width=\columnwidth]{../figs/scatter.pdf}}
\caption{(a) Leave-one-family-out predictions of the log-linear model against Monte Carlo cover times, colored by family. (b) Generic upper bounds against measured cover times on log--log axes; the dashed line is equality.}
\label{fig:scatter}
\end{center}
\vskip -0.2in
\end{figure}

\section{Discussion and Limitations}
\label{sec:disc}
The predictor is useful when the graph is an ordinary member of a family whose structure is captured by $n$, $m$, $\Delta$ and $D$, such as random, regular, scale-free, and small-world graphs at this size. It is unreliable on graphs with a pronounced bottleneck or a dense core attached to a long path, because those structures set the cover time through quantities the four statistics do not see. A richer feature set, such as the conductance or the maximum hitting time, would likely help. We did not test this because those quantities are close in cost to the Matthews bound itself.

The limitations are substantial. (i) The graphs are synthetic and their size range is small; the fit should not be extrapolated to $n\gg 30$, where the lollipop-type extremal behaviour appears. (ii) The target is a Monte Carlo estimate with a few percent of noise, and the max over start vertices is taken over sampled starts for $n>12$, so it may underestimate the true $C(G)$ by a small amount. (iii) The bootstrap intervals do not refit the model. (iv) Leave-one-family-out errors are inflated by the extrapolation effect discussed above. (v) We do not prove anything about these predictors, and the coefficients should be read as descriptive.

\section{Reproducibility}
\label{sec:repro}
The code is in \texttt{src/}: \texttt{graphs.py} (families), \texttt{sim.py} (vectorised walks, exact solver for $n\le 8$, hitting times), \texttt{run.py} (ground truth; \texttt{--walks 300 --seed 0}), and \texttt{analyze.py} (fits, CV, LOFO, bootstrap, figure). Run \texttt{python src/run.py --walks 300 --seed 0 --out data/results.csv} followed by \texttt{python src/analyze.py data/results.csv data}, using the Python version and package versions listed in \texttt{README.md}. Random seeds are fixed at 0 for graph generation, starts and walks; \texttt{data/analysis\_log.txt} holds the output reported above. The run takes about two minutes on one CPU.



\bibliography{refs}
\bibliographystyle{icml2025}

\end{document}
